%%%%
% Преамбула: подключение необходимых пакетов
% Редактируйте осторожно!
%

\documentclass[hyperref={unicode}]{beamer}

\usepackage[utf8x]{inputenc}
\usepackage[english, russian]{babel}
\usepackage{color, colortbl}
\usepackage{rotating} 
\usepackage{graphicx}
\usepackage{algorithmic}

\usetheme[nosecheader]{PetrSU-CS}


%%%%
% Преамбула: основные параметры презентации
% Отредактируйте в соответствии с комментариями
%

\title[%
    % Краткое название работы не используется в этой презентации!
    Бобры и Интернет
]{%
    % Полное название работы отображается на титульной странице
    Модели компьютерного зрения \\
    в производственных веб-системах
}

% Подзаголовком опишите тип работы:
% - Курсовая работа
% - Выпускная квалификационная работа бакалавра
% - Дипломная работа
% - Магистерская диссертация
\subtitle{Промежуточный отчет о научно-исследовательской работе}

\author[%
    % Имя и фамилия автора работы отображаются на каждом слайде в нижнем колонтитуле
    Олег Плугин
]{%
    % Имя, отчество и фамилия автора работы отображатся на титульном слайде
    Олег Антонович Плугин
}

\date[%
    % Дата защиты
    05.06.2026
]{%
    % Руководитель
    Научный руководитель: к. ф-м. н., доцент Д. Ж. Корзун 
}

\institute[%
    % Краткое название организации не используется в этой презентации
    ПетрГУ
]{%
    % Полное название организации и подразделения
    Петрозаводский государственный университет\\
    Кафедра информатики и математического обеспечения
}


%%%%
%
% Начало содержимого слайдов
%

\begin{document}

\begin{frame}
\titlepage
\end{frame}

%------------------------------------------------
\begin{frame}{Актуальность работы}
\begin{itemize}
    \item Активное внедрение компьютерного зрения в производственных и транспортных системах
    \item Необходимость обработки видеоданных в режиме, близком к реальному времени
    \item Ограниченные вычислительные ресурсы мобильных и встраиваемых устройств
    \item Актуальность оптимизации нейросетевых моделей без существенной потери точности
\end{itemize}
\end{frame}

%------------------------------------------------
\begin{frame}{Цель и задачи текущего этапа}
\textbf{Цель работы:}
\begin{itemize}
    \item Практическая реализация и оптимизация конвейера машинного обучения для системы видеоаналитики дефектов дорожного покрытия
\end{itemize}

\medskip
\textbf{Основные задачи:}
\begin{itemize}
    \item Формирование и программная диагностика обучающей выборки
    \item Обучение модели инстанс-сегментации с применением Transfer Learning
    \item Автоматизация контроля метрик качества (mAP, Loss, Confusion Matrix)
    \item Подготовка обученной модели к интеграции (оптимизация инференса)
\end{itemize}
\end{frame}

%------------------------------------------------
\begin{frame}{Анализ научной статьи}
\begin{itemize}
    \item Рассматривается задача оптимизации моделей детекции объектов
    \item Основная проблема:
    \begin{itemize}
        \item сложная архитектура модели $\rightarrow$ высокая вычислительная сложность
        \item упрощение модели $\rightarrow$ критическое снижение качества распознавания (mAP)
    \end{itemize}
    \item Цель — поиск компромисса между точностью (mAP) и производительностью (FPS)
    \item Условия эксплуатации: ограниченные ресурсы ЦП и памяти
\end{itemize}
\end{frame}

%------------------------------------------------
\begin{frame}{Параметры и метрики оптимизации}
\textbf{Параметры модели:}
\begin{itemize}
    \item количество слоёв нейросети
    \item размер входного изображения
    \item количество и размеры якорных рамок
\end{itemize}

\medskip
\textbf{Метрики оценки:}
\begin{itemize}
    \item mAP, Precision, Recall
    \item IoU
    \item FPS и время обработки кадра
\end{itemize}
\end{frame}

%------------------------------------------------
\begin{frame}{Алгоритм повышения точности и производительности}
\begin{itemize}
    \item Итеративный подбор количества слоёв
    \item Изменение размера входного изображения
    \item Оптимизация якорных рамок методом кластеризации
    \item Нормализация якорей под карту признаков
    \item Оценка метрик и выбор оптимальной конфигурации
\end{itemize}
\end{frame}

%------------------------------------------------
\begin{frame}{Результаты оптимизации модели}
\begin{itemize}
    \item Уменьшение числа слоёв нейросети (с 24 до 9)
    \item Рост метрик точности (mAP, Precision, Recall)
    \item Увеличение производительности до 5--10 раз
    \item Возможность применения на мобильных и малопроизводительных устройствах
\end{itemize}
\end{frame}

%------------------------------------------------
\begin{frame}{Система видеоаналитики дорожного покрытия}
\begin{itemize}
    \item Автоматическое обнаружение и инстанс-сегментация дефектов дорог
    \item Разработка в рамках грантового проекта
    \item Источники данных:
    \begin{itemize}
        \item видеокамеры мобильной лаборатории
        \item данные дополнительных сенсоров
    \end{itemize}
    \item Переход от теоретического анализа к практической программной реализации конвейера машинного обучения
\end{itemize}
\end{frame}

%------------------------------------------------
% НОВЫЙ СЛАЙД 1
\begin{frame}{Формирование обучающей выборки (Датасета)}
\begin{itemize}
    \item Использование платформы \textbf{CVAT} для первичной разметки видеокадров полигонами
    \item Программная диагностика аннотаций и утверждение \textbf{6 целевых классов} дефектов
    \item Разработка Python-скриптов для слияния разрозненных задач разметки:
    \begin{itemize}
        \item Объединено 15 задач
        \item Обработано более 58 000 изображений
    \end{itemize}
    \item Программная конвертация итогового датасета в стандарт \textbf{COCO} (оптимально для инстанс-сегментации)
\end{itemize}
\end{frame}

%------------------------------------------------
% НОВЫЙ СЛАЙД 2
\begin{frame}{Архитектура нейросети и процесс обучения}
\begin{itemize}
    \item Выбор одностадийной архитектуры \textbf{RTMDet-Ins} (баланс state-of-the-art точности и высокого FPS)
    \item Развертывание среды на базе \textbf{PyTorch} и \textbf{MMDetection}
    \item Применение \textbf{Transfer Learning}: 
    \begin{itemize}
        \item Инициализация весами, предобученными на датасете COCO
        \item Радикальное ускорение сходимости модели
    \end{itemize}
    \item Итеративная настройка гиперпараметров (learning rate, batch size) и применение механизмов регуляризации (weight decay)
\end{itemize}
\end{frame}

%------------------------------------------------
% НОВЫЙ СЛАЙД 3
\begin{frame}{Оценка качества модели и аналитика}
\begin{itemize}
    \item Автоматизация аналитики: разработка скриптов для парсинга логов обучения
    \item Анализ графиков \textbf{функции потерь (Loss)} для контроля сходимости и предотвращения переобучения (overfitting)
    \item Оценка пространственной точности (\textbf{IoU}) и комплексной метрики \textbf{mAP}
    \item Построение \textbf{матрицы ошибок (Confusion Matrix)}:
    \begin{itemize}
        \item Выявление ложноположительных (FP) и ложноотрицательных (FN) срабатываний
        \item Определение проблемных классов для балансировки датасета
    \end{itemize}
\end{itemize}
\end{frame}

%------------------------------------------------
% НОВЫЙ СЛАЙД 4
\begin{frame}{Оптимизация инференса и интеграция}
\begin{itemize}
    \item Проблема: избыточность и тяжеловесность исходного формата весов PyTorch для production-среды
    \item Решение: экспорт вычислительного графа в открытый стандарт \textbf{ONNX} (Open Neural Network Exchange)
    \item Преимущества подхода:
    \begin{itemize}
        \item Независимость инференса от исходных фреймворков
        \item Минимизация потребления ОЗУ и видеопамяти
        \item Ускорение обработки кадров за счет \textbf{ONNX Runtime}
    \end{itemize}
    \item Подготовка базы для интеграции модуля в веб-систему
\end{itemize}
\end{frame}

%------------------------------------------------
\begin{frame}{Итоги и дальнейшая работа}
\textbf{Достигнутые результаты текущего этапа:}
\begin{itemize}
    \item Сформирован унифицированный датасет (COCO)
    \item Успешно обучена модель инстанс-сегментации RTMDet-Ins
    \item Внедрен автоматизированный мониторинг метрик
\end{itemize}

\medskip
\textbf{Направления дальнейших исследований:}
\begin{itemize}
    \item Практическая реализация скриптов конвертации весов в ONNX
    \item Интеграция оптимизированного вычислительного модуля в разрабатываемую веб-систему
    \item Проведение нагрузочных тестов на новых видеоданных
\end{itemize}
\end{frame}

\begin{frame}
  \frametitle{}
  
{\Large\mbox{}\hfil Спасибо за внимание!}
  
\end{frame}
\end{document}
